A biological cell is a conglomerate of many interacting biochemical components organised in sub-units. Recently, it has been established that a cell employs phase separation \cite{alberti_phase_2017} to organise itself spatially and form membrane-less compartments (see Figure \ref{fig:phase_sep}). These membrane-less compartments can adapt to external stress, help regulate genes or even cause some neurodegenerative diseases \cite{dolgin_what_2018, hyman_liquid-liquid_2014}. Hence, it becomes increasingly important to study under what conditions such membrane-less compartments form and function. Estimating their free energies serves as a good starting point.\\

\begin{figure}[h!]	\includegraphics[width=0.25\linewidth]{imgs/intro/phase_sep.png}
	\caption{Phase separation in nucleoli of X. \textit{laevis}. Adapted from \cite{feric_coexisting_2016, dolgin_how_2024}.}
	\label{fig:phase_sep}
\end{figure}

The Flory-Huggins formulation can be used to study the free energies of phase-separating materials. This formulation is often used to study systems with two components \cite{rubinstein_polymer_2003}. Extending it to multiple interacting components is straightforward and allows one to model a cell as numerous phase-separated compartments. However, the dimensionality of the corresponding free energy landscape scales accordingly. Several methods have been developed to study phase coexistence but are still limited to a few components \cite{sch2023, mao_phase_2019}.\\

This report explores methods to sample the free energy landscape of a system with $n$ components and $m$-distinct compartments, whereby a compartment is a spatial domain with constant concentrations of each component. We also propose a general method for systematically reducing the landscape's dimensionality. We illustrate the method in the light of elementary systems but with considerably high dimensionality of the free energy landscape..\\

The remainder of this report is structured as follows: Section 2 briefly outlines the underlying theory of Flory-Huggins' free energy formulation and its extension to multi-phase systems. Following this, the approach of sampling free-energy landscapes with a simple random walk is introduced. We then describe our approach to reduce the dimensions of an $n$-component system. Section 3 shows the results from our methods for particular example systems. Lastly, section 4 discusses further steps in light of the results obtained.